\section{Introduction}
\label{sec:introduction}

Large language models (LLMs) are moving from single-question assistants to
components of expert and decision-support systems that keep a task alive over
many turns. A procurement assistant tracks a budget as invoices arrive, a
logistics assistant follows stock across warehouses, and a laboratory assistant
records volumes as reagents are transferred. In each case the user supplies
information piece by piece, corrects earlier statements and asks questions
whose answers depend on everything said so far. The model's context therefore
accumulates two things: its own earlier replies and a task state that changes
at every turn. We call this accumulated context the \emph{conversation
history}. In our experiments, the condition that retains it is called
\emph{dialogue history}.

There is growing evidence that this setting is costly.
\citet{laban2025lost} showed that performance drops substantially when the
information of a single-turn task is revealed over several turns: models commit
to early answers and fail to recover. \citet{huang2026benefit} showed that
removing the assistant's own previous replies often preserves or improves
response quality, a harm they call \emph{context pollution}. Long-context
benchmarks show that information in long contexts is not used uniformly
\citep{liu2024lost}, while prior assistant errors can propagate through later
turns \citep{huang2026benefit}.

For a system designer, two operational questions follow, and prior work does
not settle either of them. The first is whether degradation persists, or
accelerates, over long horizons once task difficulty is controlled. Long
dialogues naturally change what they ask, so a falling accuracy curve may
reflect harder questions rather than a harder context. The second is whether
the degradation can be detected \emph{before} accuracy fails. A deployed system
rarely observes the correct answer, so it must rely on observable process
signals. Reasoning quality has been decomposed into dimensions such as
agreement between samples, robustness to rewording, logical coherence of the
reasoning chain and stability of the trace \citep{rqeval2026}. If these
dimensions deteriorated before accuracy, they could serve as an early-warning
layer, and classic drift detectors could turn them into alarms with controlled
false-alarm rates \citep{gama2014survey,lu2019learning}. ERGO uses token-entropy
spikes in a related role \citep{khalid2025ergo}. To our knowledge, whether
established process metrics warn under accumulated history has not been tested
under prospectively fixed analyses with calibrated false alarms.

We address both questions with two complementary studies. Each was analysed
under a protocol frozen before its data were collected, and every analysis
added later is labelled as exploratory or post-hoc.
\emph{Study~1} uses the math tasks of Lost in Conversation
\citep{laban2025lost}. At every turn, each model answers twice from identical
user information, once with and once without its own earlier replies, which
yields a task-paired estimate of the net effect of retaining assistant-side
history. \emph{Study~2} introduces CDR-Stream, a purpose-built benchmark of
100-turn state-tracking dialogues with exact per-turn answers, difficulty
calibrated separately for each model, and scheduled difficulty shifts. Its
control arm replaces the history with a solver-verified state summary. Across
both studies we then test the measurable quality metrics, a paired readiness
measure and window-based drift detectors as early-warning signals. The two
studies use different tasks, models and controls. We therefore read them as
revealing complementary failure modes, not as one dose--response curve.

The paper makes four contributions.
\begin{enumerate}
  \item A frozen, task-paired confirmatory estimate of the effect of retaining
    assistant replies on exactly scored multi-turn math tasks. It replicates
    the direction reported by \citet{huang2026benefit} under a prospectively
    fixed analysis; it is a replication, not a discovery.
  \item CDR-Stream, a long-horizon benchmark with exact per-turn gold and
    per-model difficulty calibration. In it, accuracy under full history
    collapsed within about twenty turns, even at matched difficulty, while a
    verified state summary preserved it in all three models tested.
  \item A prospective test of early warning. The measurable consistency,
    coherence and stability metrics did not predict failure. Coherence and
    stability even stayed level or rose as accuracy fell. A paired readiness
    measure reduced to the persistence of the current unresolved state, and
    single-stream detectors were too weak for the observed drifts.
  \item Design guidance for LLM-based expert systems, supported by a detector
    power simulation and a detector audit: manage state explicitly, monitor by
    comparison or aggregation across sessions, and verify that a baseline
    detector can fire on streams of the intended length.
\end{enumerate}

Section~\ref{sec:related} reviews related work and positions the monitoring
framework. Sections~\ref{sec:study1} and~\ref{sec:study2} present the studies,
their metrics and evidence standards, and
Section~\ref{sec:monitoring} asks whether the degradation can be detected
early. Section~\ref{sec:discussion} discusses implications and limitations.
